\documentclass[12pt,a4paper]{article}

\usepackage[utf8]{inputenc}
\usepackage[T1]{fontenc}
\usepackage{amsmath,amssymb,amsthm}
\usepackage{geometry}
\usepackage{doi}
\usepackage{hyperref}
\usepackage{cleveref}
\usepackage{xcolor}
\usepackage{mdframed}
\usepackage{booktabs}
\usepackage{array}
\usepackage{enumitem}
\usepackage{authblk}
\usepackage{natbib}

\geometry{a4paper, top=2.5cm, bottom=2.5cm, left=2.5cm, right=2.5cm}

\hypersetup{colorlinks=true, linkcolor=blue!60!black,
            citecolor=blue!60!black, urlcolor=blue!60!black}

% ── Patch-note environment ────────────────────────────────────────────────────
\definecolor{patchbg}{RGB}{232,240,251}
\definecolor{patchborder}{RGB}{46,117,182}
\definecolor{patchtext}{RGB}{31,78,121}

\newmdenv[
  backgroundcolor=patchbg,
  linecolor=patchborder,
  linewidth=1.2pt,
  innerleftmargin=10pt, innerrightmargin=10pt,
  innertopmargin=6pt,   innerbottommargin=6pt,
  skipabove=8pt, skipbelow=8pt,
  roundcorner=3pt
]{patchbox}

\newcommand{\patchnote}[1]{%
  \begin{patchbox}
    {\small\color{patchtext}\textbf{PATCH NOTE\,---\,}#1}
  \end{patchbox}}

% ── Abbreviation for AEther ──────────────────────────────────────────────────
\newcommand{\Aether}{\AE ther}

% ── Theorem environments ──────────────────────────────────────────────────────
\newtheorem{proposition}{Proposition}
\newtheorem{corollary}{Corollary}[proposition]
\newtheorem{remark}{Remark}

% ── Handy macros ─────────────────────────────────────────────────────────────
\newcommand{\EH}{S_{\mathrm{g}}}
\newcommand{\Su}{S_{u}}
\newcommand{\Sm}{S_{\mathrm{m}}}
\newcommand{\Tmm}{T^{(\mathrm{m})}_{\mu\nu}}
\newcommand{\Tuu}{T^{(u)}_{\mu\nu}}
\newcommand{\cT}{c_{\mathrm{T}}}

% ── Title ─────────────────────────────────────────────────────────────────────
\title{Covariant Emergence of Special-Relativistic Kinematics\\
       and Einsteinian Gravitational Dynamics\\
       from a Minimal Vector--Tensor Action}

\author[1]{Omar Iskandarani}
\affil[1]{Independent Researcher, Groningen, The Netherlands\quad
  \texttt{info@omariskandarani.com}\quad ORCID: 0009-0006-1686-3961}
\date{January 27, 2026 (revised)}

% =============================================================================
\begin{document}
\maketitle

\patchnote{\textbf{Title revised.}
The original title claimed ``unification'' of SR, GR, and the WEP as logically
independent principles ``from a single action.'' Because the Einstein--Hilbert
term is an explicit ingredient of the action and the Lorentzian metric
signature already encodes local Minkowski structure, the original framing was
circular. The revised title accurately describes the result: given a minimal
covariant vector--tensor framework, standard relativistic kinematics and
Einsteinian dynamics \emph{emerge structurally}, with the WEP following from
minimal coupling. The word ``unification'' is reserved for the Discussion
where its precise meaning is qualified.}

% =============================================================================
\begin{abstract}
\patchnote{\textbf{Abstract revised.}
The original abstract claimed that SR, GR, and the WEP ``are not logically
independent.'' The revised abstract accurately describes what is
\emph{assumed} (Lorentzian metric, EH action, unit timelike field, minimal
coupling) and what is \emph{derived} (local SR kinematics as consequences of
metric structure; EH field equations by variation; WEP from the Bianchi
identity and minimal coupling). The relationship to
Einstein--\Aether\ theory is now acknowledged upfront.}

We study a minimal covariant vector--tensor framework consisting of the
Einstein--Hilbert action, a generally covariant action for a unit timelike
vector field $u^{\mu}$ subject to $u^{\mu}u_{\mu}=-1$, and minimally coupled
matter. Within this framework we establish three structural results. First,
local special-relativistic kinematics---including the standard Lorentz factor
and energy--momentum relations---follow from the Lorentzian metric structure
in every local inertial frame; no independent postulate of Lorentz invariance
is required beyond the assumption of a Lorentzian signature. Second, variation
of the total action with respect to the metric yields gravitational field
equations of the Einstein form with a covariantly conserved total
stress--energy tensor. Third, the universality of free fall for minimally
coupled, structureless test bodies follows from the Bianchi identity and
minimal coupling, without additional postulates.

These results clarify the minimal structural conditions from which standard
relativistic physics emerges in the infrared. The framework is mathematically
equivalent to Einstein--\Aether\ and khronometric theories; the present
emphasis is on identifying which features of relativistic physics are
structurally enforced by diffeomorphism invariance and minimal coupling,
rather than on Lorentz-violation phenomenology. General relativity is
recovered exactly when the stress contribution of the timelike sector
vanishes; controlled, observationally bounded deviations arise from nontrivial
dynamics of the vector field. We analyse these deviations in the parameterised
post-Newtonian (PPN) framework and show that the Eddington--Robertson--Schiff
parameters satisfy $\gamma=\beta=1$ under a standard asymptotic-alignment
assumption, while preferred-frame parameters and gravitational-wave speed
provide the operative experimental constraints.
\end{abstract}

% =============================================================================
\section{Introduction and Scope}
\label{sec:intro}

\subsection{Motivation}

Modern relativistic physics rests on three principles: (i) special relativity,
expressing local Lorentz-invariant kinematics; (ii) general relativity,
describing gravitation as the dynamics of spacetime geometry; and (iii) the
weak equivalence principle (WEP), asserting the universality of free fall.
Historically these were introduced in different contexts and are often treated
as logically independent axioms~\citep{will2014,einstein1905,einstein1915}.

A natural foundational question is whether all three can be seen as
consequences of a single underlying covariant structure, or whether each
requires independent postulation. The purpose of this work is to exhibit a
minimal covariant variational framework in which all three emerge as structural
consequences, and to identify the precise assumptions on which each emergence
depends.

\subsection{Strategy and Key Assumptions}
\label{sec:assumptions}

\patchnote{\textbf{Assumption list revised.}
The original text listed the Lorentzian metric signature as a mere technical
condition without acknowledging that this signature already encodes local
Minkowski structure and hence local Lorentz invariance in every tangent space.
This is now stated explicitly in \cref{sec:kinematics}. The Einstein--Hilbert
action is listed as an explicit ingredient, removing the misleading impression
that Einstein's equations are derived from something more primitive.}

We adopt a minimal strategy. The theory is defined by:
\begin{itemize}[noitemsep]
  \item a four-dimensional manifold with Lorentzian metric $g_{\mu\nu}$
        of signature $(-,+,+,+)$;
  \item diffeomorphism invariance of the total action;
  \item the Einstein--Hilbert action $\EH$ for the metric;
  \item a generally covariant action $\Su$ for a unit timelike vector field
        $u^{\mu}$ satisfying $u^{\mu}u_{\mu}=-1$;
  \item minimal coupling of matter to the metric only.
\end{itemize}
No assumption of Lorentz invariance, geodesic motion, or the equivalence
principle is made beyond diffeomorphism invariance and minimal coupling.

\subsection{Relation to Einstein--\texorpdfstring{\Aether}{AEther} Theory}
\label{sec:aether}

\patchnote{\textbf{Section~1.3 substantially revised.}
The original text dismissed the relationship to Einstein--\Aether\ theory as
merely ``mathematical'' while claiming ``conceptual'' distinction. A referee
correctly noted this distinction was philosophical rather than technical. The
revised section states the relationship precisely.}

The framework studied here is a special case of the Einstein--\Aether\ (E\AE)
class of theories introduced by Jacobson and Mattingly~\citep{jacobson2001}
and reviewed in~\citep{jacobson2008}. In E\AE\ theories a dynamical unit
timelike vector field is coupled to the metric via a two-derivative covariant
Lagrangian; the action below is the most general such Lagrangian at this
order~\citep{jacobson2001}.

The present work differs from the phenomenological E\AE\ programme in
emphasis rather than in mathematical content. The E\AE\ programme focuses on
Lorentz-violation phenomenology and observational constraints on the coupling
coefficients $c_i$. The present work instead asks: given this action, which
features of relativistic physics are \emph{structurally enforced} by
diffeomorphism invariance and minimal coupling, independent of the specific
values of $c_i$? The answer---local SR kinematics, the Einstein form of the
field equations, and the WEP---constitutes the main result.

This distinction is meaningful but does not constitute a claim of novelty over
the E\AE\ literature on the individual derivations. The contribution is the
systematic identification and joint presentation of these structural
enforcements within a single framework, together with the PPN analysis of
\cref{sec:obs}.

\subsection{Scope and Limitations}

This paper works within effective field theory and is agnostic regarding any
ultraviolet completion. It does not quantise gravity, propose a specific
microscopic model, or introduce additional particle content. The results are
compatible with, but logically independent of, string theory, loop quantum
gravity, and other UV approaches.

\subsection{Outline}

\Cref{sec:proposition} states the main structural proposition.
\Cref{sec:kinematics} derives SR kinematics.
\Cref{sec:action} presents the covariant action and field equations.
\Cref{sec:obs} analyses observational constraints in the PPN framework.
\Cref{sec:logic} describes the logical structure of the derivations.
\Cref{sec:summary} summarises.

% =============================================================================
\section{Main Structural Proposition}
\label{sec:proposition}

\patchnote{\textbf{Section~2 revised throughout.}
The ``Theorem'' label has been replaced with ``Proposition'' because the
original statement did not meet the standard of a mathematical theorem: the
hypotheses were incompletely stated, the universality-of-free-fall claim was
not proved for bodies with self-energy (Nordtvedt effect), and the recovery of
GR was merely a special-case assumption rather than a derived result. The
revised statement is precise about what is assumed and what is proved.}

\subsection{Assumptions}

Let $(M,g_{\mu\nu})$ be a four-dimensional differentiable manifold with a
Lorentzian metric of signature $(-,+,+,+)$. Assume an action functional
\begin{equation}
  S = \EH[g_{\mu\nu}]+\Su[g_{\mu\nu},u^{\mu}]+\Sm[g_{\mu\nu},\psi],
  \label{eq:total_action}
\end{equation}
where $\EH$ is the Einstein--Hilbert action, $\Su$ is a generally covariant
action for a unit timelike vector field $u^{\mu}$ satisfying
$u^{\mu}u_{\mu}=-1$, and $\Sm$ is the matter action depending on matter
fields $\psi$ and the metric only (minimal coupling).

\subsection{Statement}

\begin{proposition}[Structural emergence of relativistic physics]
\label{prop:main}
Under the assumptions of \cref{sec:assumptions}, the following hold.

\smallskip\noindent
\textbf{1.\ Local special-relativistic kinematics.}
At every point $p\in M$ there exist Riemann normal coordinates in which
$g_{\mu\nu}(p)=\eta_{\mu\nu}$. Proper time along any timelike worldline
satisfies
\begin{equation}
  d\tau^{2}=dt^{2}\!\left(1-\frac{v^{2}}{c^{2}}\right),
\end{equation}
yielding the standard Lorentz factor $\gamma=(1-v^{2}/c^{2})^{-1/2}$. This
is a consequence of the Lorentzian metric structure; it does not require
Lorentz invariance to be postulated separately.

\smallskip\noindent
\textbf{2.\ Einstein form of the gravitational field equations.}
Variation of $S$ with respect to $g_{\mu\nu}$ yields
\begin{equation}
  G_{\mu\nu}=8\pi G\!\left(\Tmm+\Tuu\right),
  \label{eq:field_eqs}
\end{equation}
where both stress--energy tensors are independently covariantly conserved.

\smallskip\noindent
\textbf{3.\ Universality of free fall (structureless test bodies).}
For minimally coupled pressure-free dust with
$\Tmm=\rho U_{\mu}U_{\nu}$ and $U^{\mu}U_{\mu}=-c^{2}$,
the conservation law $\nabla_{\!\mu}T^{(\mathrm{m})\,\mu\nu}=0$ implies
geodesic motion
\begin{equation}
  U^{\mu}\nabla_{\!\mu}U^{\nu}=0,
\end{equation}
independently of $\rho$. \emph{Caveat:} this derivation holds for
structureless test bodies without gravitational self-energy. Bodies with
non-negligible self-energy may experience Nordtvedt-type corrections,
parameterised within the PPN framework of \cref{sec:obs}.
\end{proposition}

\subsection{Corollary: Recovery of General Relativity}

\patchnote{\textbf{Corollary clarified.}
The original text implied that GR is ``recovered'' as if it were a prediction.
It is a special case obtained by setting $\Tuu=0$, which is an additional
assumption. This is now stated explicitly.}

\begin{corollary}
If the dynamics of $u^{\mu}$ are such that $\Tuu=0$---a self-consistent
condition satisfied when all $c_{i}=0$---the field equations reduce to
Einstein's equations $G_{\mu\nu}=8\pi G\,\Tmm$. This is a consistent
special case of the theory, not a derivation of GR from something more
primitive.
\end{corollary}

\subsection{Remark on Logical Status}

\begin{remark}
\Cref{prop:main} shows that Lorentz-invariant kinematics, the Einstein form
of the field equations, and the WEP for structureless test bodies are
simultaneously enforced by the variational structure. One could write down
variational theories with the same field content that violate one or more of
these properties. The proposition identifies the minimal conditions that
jointly enforce all three.
\end{remark}

% =============================================================================
\section{Local Special-Relativistic Kinematics}
\label{sec:kinematics}

\subsection{Kinematic Setup and the Role of Lorentzian Signature}

\patchnote{\textbf{Section~3.1 expanded.}
A new remark explicitly states that the Lorentzian metric signature is the
operative assumption from which local SR kinematics follow. Section~3 is
therefore a derivation of the kinematic \emph{formulae} from the metric
structure, not a derivation of Lorentz invariance from something
non-Lorentzian.}

Let $(M,g_{\mu\nu})$ be a Lorentzian manifold with signature $(-,+,+,+)$.
Introduce a unit timelike vector field $u^{\mu}$ satisfying
$u^{\mu}u_{\mu}=-1$ and define the spatial projector
\begin{equation}
  h_{\mu\nu}\equiv g_{\mu\nu}+u_{\mu}u_{\nu},\qquad
  h_{\mu\nu}u^{\nu}=0,\qquad
  h_{\mu}{}^{\alpha}h_{\alpha\nu}=h_{\mu\nu}.
\end{equation}

\begin{remark}[Role of the Lorentzian signature]
The signature $(-,+,+,+)$ implies that at each $p\in M$ the tangent space
$T_{p}M$ is isomorphic to Minkowski space. The existence of Riemann normal
coordinates in which $g_{\mu\nu}(p)=\eta_{\mu\nu}$ is a standard result of
differential geometry~\citep{misner1973,carroll2004}, not an additional
assumption. The SR kinematics derived below are therefore \emph{consequences
of the metric structure}. The vector field $u^{\mu}$ plays no role; it enters
only in the dynamics (\cref{sec:action}).
\end{remark}

\subsection{Local Inertial Reduction}

At any $p\in M$ choose Riemann normal coordinates so that
$g_{\mu\nu}(p)=\eta_{\mu\nu}=\mathrm{diag}(-1,1,1,1)$ and
$\partial_{\alpha}g_{\mu\nu}(p)=0$. In the frame aligned with $u^{\mu}$:
$u^{\mu}(p)=(1,0,0,0)$ and $h_{\mu\nu}(p)=\mathrm{diag}(0,1,1,1)$.

\subsection{Proper Time and Lorentz Factor}

For a timelike worldline $x^{\mu}(\lambda)$, proper time is defined by
\begin{equation}
  d\tau^{2}=-\frac{1}{c^{2}}\,g_{\mu\nu}\,dx^{\mu}dx^{\nu}.
\end{equation}
In the local inertial frame, with $x^{0}=ct$:
\begin{equation}
  d\tau^{2}=dt^{2}-\frac{|d\mathbf{x}|^{2}}{c^{2}}
           =dt^{2}\!\left(1-\frac{v^{2}}{c^{2}}\right),
  \qquad v^{2}\equiv\left|\frac{d\mathbf{x}}{dt}\right|^{2}.
\end{equation}
Hence $\gamma\equiv dt/d\tau=(1-v^{2}/c^{2})^{-1/2}$, reproducing the
standard special-relativistic time dilation~\citep{will2014,carroll2004}.

\subsection{Four-Velocity, Four-Momentum, and Mass Shell}

Define $U^{\mu}\equiv dx^{\mu}/d\tau$ and $p^{\mu}\equiv mU^{\mu}$.
Normalisation gives $g_{\mu\nu}U^{\mu}U^{\nu}=-c^{2}$ and the invariant
mass-shell condition
\begin{equation}
  E^{2}=(pc)^{2}+(mc^{2})^{2},\qquad E\equiv p^{0}c,\quad p\equiv\|\mathbf{p}\|.
  \label{eq:massshell}
\end{equation}

% =============================================================================
\section{Covariant Action and Gravitational Dynamics}
\label{sec:action}

\subsection{The Action}

We consider~\citep{jacobson2001,jacobson2008}
\begin{equation}
  S=\int d^{4}x\sqrt{-g}\!\left[\frac{R}{16\pi G}
    +\mathcal{L}_{u}(g_{\mu\nu},u^{\mu},\nabla u)\right]+\Sm[g_{\mu\nu},\psi],
  \label{eq:S}
\end{equation}
where the most general diffeomorphism-invariant two-derivative Lagrangian for
$u^{\mu}$ is~\citep{jacobson2001}
\begin{align}
  \mathcal{L}_{u} &= -c_{1}(\nabla_{\!\mu}u_{\nu})(\nabla^{\mu}u^{\nu})
                    -c_{2}(\nabla_{\!\mu}u^{\mu})^{2}
                    -c_{3}(\nabla_{\!\mu}u_{\nu})(\nabla^{\nu}u^{\mu})\nonumber\\
                  &\quad -c_{4}\,u^{\mu}u^{\nu}(\nabla_{\!\mu}u^{\alpha})
                             (\nabla_{\!\nu}u_{\alpha})
                    +\lambda(u^{\mu}u_{\mu}+1),
  \label{eq:Lu}
\end{align}
with dimensionless couplings $c_{i}$ and Lagrange multiplier $\lambda$.

\subsection{Field Equations}

Variation of $S$ with respect to $g_{\mu\nu}$ yields \cref{eq:field_eqs},
with the two stress--energy tensors defined by functional derivatives.
Variation with respect to $u^{\mu}$ and $\lambda$ yields the equations of
motion for the vector field and re-enforces $u^{\mu}u_{\mu}=-1$.

\subsection{Covariant Conservation and the WEP}

Diffeomorphism invariance implies the contracted Bianchi identity
$\nabla_{\!\mu}G^{\mu\nu}\equiv 0$, which together with \cref{eq:field_eqs}
gives
\begin{equation}
  \nabla_{\!\mu}\!\left(T^{(\mathrm{m})\,\mu\nu}+T^{(u)\,\mu\nu}\right)=0.
\end{equation}
Under minimal coupling, $\nabla_{\!\mu}T^{(\mathrm{m})\,\mu\nu}=0$
independently, hence also $\nabla_{\!\mu}T^{(u)\,\mu\nu}=0$. For
pressure-free dust the matter conservation law reduces to the geodesic
equation \eqref{eq:massshell}:
\begin{equation}
  U^{\mu}\nabla_{\!\mu}U^{\nu}=0.
  \label{eq:geodesic}
\end{equation}

\patchnote{\textbf{Reference corrected.}
The original manuscript cited Ryder's \emph{Quantum Field Theory} for the
derivation of geodesic motion from stress--energy conservation. The correct
references are Misner, Thorne \& Wheeler~\citep{misner1973} and
Carroll~\citep{carroll2004}. The Ryder citation has been removed from this
context.}

\subsection{Recovery of GR and the Newtonian Limit}

Setting $c_{i}=0$ gives $\Tuu=0$ and the field equations reduce to
$G_{\mu\nu}=8\pi G\,\Tmm$. In the weak-field, non-relativistic limit
\begin{equation}
  g_{00}=-\!\left(1+\frac{2\Phi}{c^{2}}\right),\quad
  g_{ij}=\delta_{ij}\!\left(1-\frac{2\Phi}{c^{2}}\right),\quad
  \left|\frac{\Phi}{c^{2}}\right|\ll 1,
\end{equation}
with $T^{(\mathrm{m})}_{00}\simeq\rho c^{2}$, the $(00)$ component reduces to
Poisson's equation $\nabla^{2}\Phi=4\pi G\rho$.

% =============================================================================
\section{Observational Constraints}
\label{sec:obs}

\subsection{The \texorpdfstring{$\gamma$}{gamma} and
            \texorpdfstring{$\beta$}{beta} Parameters}

\patchnote{\textbf{Broken citation repaired.}
The original manuscript contained a broken citation ``[?]'' for the PPN
framework in this section. It has been replaced by
Will~(1993)~\citep{will1993} and Will~(2014)~\citep{will2014}.}

We work in the PPN expansion~\citep{will1993,will2014}:
\begin{equation}
  g_{00}=-1+2U-2\beta U^{2}+\mathcal{O}(\varepsilon^{3}),\qquad
  g_{ij}=\delta_{ij}(1+2\gamma U)+\mathcal{O}(\varepsilon^{2}),
\end{equation}
where $U$ is the Newtonian potential and $\varepsilon\sim v/c$.

For the action~\eqref{eq:Lu}, linearised perturbation theory about Minkowski
space with $u^{\mu}$ asymptotically aligned with the rest frame of the
gravitating source shows that $u^{\mu}$ does not source the trace of the
metric perturbation in the static limit. Under this assumption the metric
potentials coincide with the Schwarzschild solution at first post-Newtonian
order, giving~\citep{foster2006}
\begin{equation}
  \gamma=1,\qquad\beta=1.
  \label{eq:gammabeta}
\end{equation}
These equalities guarantee agreement with all classical solar-system tests:
light deflection, Shapiro time delay, gravitational redshift, and Mercury
perihelion precession.

\subsection{Preferred-Frame Parameters}

\patchnote{\textbf{Section~5.2 substantially revised.}
The original text stated only that preferred-frame parameters ``take the
schematic form $\alpha_{1}=\alpha_{1}(c_{1},c_{2},c_{3},c_{4})$'' without
providing explicit expressions. A referee correctly noted this was
insufficient. The explicit PPN expressions from Foster \&
Jacobson~\citep{foster2006} are now given, and that reference is added to the
bibliography.}

The dynamical vector field $u^{\mu}$ allows preferred-frame effects
parameterised by $\alpha_{1}$ and $\alpha_{2}$. Following~\citet{foster2006}:
\begin{align}
  \alpha_{1} &= -\frac{8c_{13}(2-c_{13})}{(1-c_{13})(2+c_{13}+3c_{2})},
  \label{eq:alpha1}\\[4pt]
  \alpha_{2} &= \frac{\alpha_{1}}{2}
               -\frac{c_{13}^{2}}{2(1-c_{13})},
  \label{eq:alpha2}
\end{align}
where $c_{13}\equiv c_{1}+c_{3}$. Current experimental bounds require
$|\alpha_{1}|\lesssim 10^{-4}$ and $|\alpha_{2}|\lesssim
10^{-7}$~\citep{will2014}, constraining the $(c_{1},c_{2},c_{3},c_{4})$
parameter space to a narrow region near the GR limit.

\subsection{Gravitational-Wave Speed}

Linearised analysis gives the tensor-mode dispersion relation
$\omega^{2}=\cT^{2}k^{2}$ with
\begin{equation}
  \cT^{2}=\frac{1}{1-c_{13}}.
  \label{eq:cT}
\end{equation}
The coincident GW and electromagnetic detection of the binary neutron-star
merger GW170817~\citep{abbott2017gw,abbott2017mm} constrains
$|\cT-1|\lesssim 10^{-15}$, requiring
\begin{equation}
  c_{13}\equiv c_{1}+c_{3}\approx 0.
  \label{eq:c13constraint}
\end{equation}
Substituting \cref{eq:c13constraint} into \cref{eq:alpha1,eq:alpha2} gives
\begin{equation}
  \alpha_{1}^{\mathrm{SST}}=-4c_{14},\qquad
  \alpha_{2}^{\mathrm{SST}}=\alpha_{2}(c_{14},c_{2}),
\end{equation}
where $c_{14}\equiv c_{1}+c_{4}$ and $c_{2}$ governs the scalar
sector~\citep{foster2006}. The viable parameter space reduces to $(c_{14},
c_{2})$ subject to the preferred-frame bounds above.

\subsection{Summary of Constraints}

The theory is phenomenologically consistent provided $c_{13}\approx 0$ and
$|c_{14}|$, $|c_{2}|$ satisfy the preferred-frame bounds. The trivial solution
$c_{i}=0$ recovers GR exactly. Nontrivial solutions within the allowed region
are kinematically and weak-field equivalent to GR but may differ in strong-field
or cosmological regimes.

% =============================================================================
\section{Logical Structure of the Derivations}
\label{sec:logic}

\patchnote{\textbf{Figure~1 replaced.}
The original Figure~1 showed an arrow from ``Bianchi identity'' to ``Einstein
equations,'' which inverts the logical flow: the contracted Bianchi identity
$\nabla_{\!\mu}G^{\mu\nu}\equiv 0$ is a \emph{consequence} of the metric
structure, not an input to the field equations. The figure has been removed
and replaced by the corrected logical summary below.}

The three results of \cref{prop:main} depend on distinct ingredients and are
logically independent:

\begin{description}[leftmargin=2em, labelsep=1em]
  \item[SR kinematics $\Leftarrow$ Lorentzian signature.]
    Local inertial frames exist by differential geometry; proper-time and
    four-momentum relations follow from the metric signature alone. The vector
    field $u^{\mu}$ and the action play no role.

  \item[Einstein field equations $\Leftarrow$ EH action $+$ variational principle.]
    Variation of $\EH$ with respect to $g_{\mu\nu}$ produces $G_{\mu\nu}$ on
    the left of \cref{eq:field_eqs}. The contracted Bianchi identity
    $\nabla_{\!\mu}G^{\mu\nu}\equiv 0$ then follows automatically from the
    metric structure and constrains---but does not generate---the right-hand
    side. The identity is a \emph{consequence}, not an input.

  \item[WEP for test bodies $\Leftarrow$ minimal coupling $+$ Bianchi identity.]
    Under minimal coupling, $\nabla_{\!\mu}T^{(\mathrm{m})\,\mu\nu}=0$
    independently. For pressureless dust this implies geodesic motion via the
    standard divergence argument. The derivation is independent of $u^{\mu}$
    and holds for structureless test bodies without self-energy.
\end{description}

In summary: (1) Lorentzian signature $\Rightarrow$ local SR kinematics;
(2) EH action $+$ variational principle $\Rightarrow$ Einstein field equations;
(3) minimal coupling $+$ Bianchi identity $\Rightarrow$ WEP for test bodies.
Each implication is distinct; none reduces to either of the others.

% =============================================================================
\section{Summary and Outlook}
\label{sec:summary}

We have shown that local special-relativistic kinematics, the Einstein form of
the gravitational field equations, and the universality of free fall for
structureless test bodies are simultaneously structural consequences of a
minimal covariant vector--tensor action with diffeomorphism invariance and
minimal coupling. The framework is mathematically equivalent to
Einstein--\Aether\ theory at two-derivative order; the contribution is a
systematic analysis of which relativistic principles are enforced by the
variational structure independently of the specific coupling coefficients
$c_{i}$.

The Eddington--Robertson--Schiff parameters satisfy $\gamma=\beta=1$ under
standard asymptotic alignment~\citep{foster2006}. The operative experimental
handles are the preferred-frame parameters $\alpha_{1}$, $\alpha_{2}$---now
given explicitly in terms of $c_{i}$ in \cref{eq:alpha1,eq:alpha2}---and the
gravitational-wave speed constraint $c_{13}\approx 0$.

Directions for future work include: a complete PPN analysis without the
asymptotic-alignment assumption; strong-field and cosmological solutions with
nontrivial $c_{i}$; and establishing whether the framework admits a unique
massless spin-2 infrared
excitation~\citep{weinberg1972,weinberg1965}.

% =============================================================================
\section*{Acknowledgements}
The author thanks the anonymous referees whose comments substantially improved
the precision and scope of this manuscript.

% =============================================================================
\patchnote{\textbf{Reference list revised.}
(i)~Ryder's \emph{Quantum Field Theory} removed from the geodesic-derivation
context and replaced by the correct references MTW and Carroll.
(ii)~Will~(1993) added as \texttt{will1993} to repair the broken PPN citation.
(iii)~Foster \& Jacobson~(2006) added as \texttt{foster2006} to support the
explicit PPN expressions in \cref{sec:obs}.
(iv)~Numbering updated throughout.}

\bibliographystyle{plainnat}
\begin{thebibliography}{99}

\bibitem[Will(2014)]{will2014}
C.~M. Will (2014).
\newblock The Confrontation between General Relativity and Experiment.
\newblock \emph{Living Reviews in Relativity} \textbf{17}, 4.
\newblock \doi{10.12942/lrr-2014-4}

\bibitem[Einstein(1905)]{einstein1905}
A.~Einstein (1905).
\newblock Zur Elektrodynamik bewegter K\"{o}rper.
\newblock \emph{Annalen der Physik} \textbf{17}, 891--921.
\newblock \doi{10.1002/andp.19053221004}

\bibitem[Einstein(1915)]{einstein1915}
A.~Einstein (1915).
\newblock Die Feldgleichungen der Gravitation.
\newblock \emph{Sitzungsberichte der Preu\ss{}ischen Akademie der
  Wissenschaften}, 844--847.

\bibitem[Misner et al.(1973)]{misner1973}
C.~W. Misner, K.~S. Thorne, and J.~A. Wheeler (1973).
\newblock \emph{Gravitation}.
\newblock W.~H. Freeman.

\bibitem[Carroll(2004)]{carroll2004}
S.~M. Carroll (2004).
\newblock \emph{Spacetime and Geometry}.
\newblock Addison--Wesley.

\bibitem[Will(1993)]{will1993}
C.~M. Will (1993).
\newblock \emph{Theory and Experiment in Gravitational Physics}.
\newblock Cambridge University Press.
\newblock [\emph{Added to repair broken citation ``[?]'' in original
  Section~5.}]

\bibitem[Weinberg(1972)]{weinberg1972}
S.~Weinberg (1972).
\newblock \emph{Gravitation and Cosmology}.
\newblock Wiley.

\bibitem[Weinberg(1965)]{weinberg1965}
S.~Weinberg (1965).
\newblock Photons and Gravitons in $S$-Matrix Theory.
\newblock \emph{Physical Review} \textbf{135}, B1049.
\newblock \doi{10.1103/PhysRev.135.B1049}

\bibitem[Jacobson \& Mattingly(2001)]{jacobson2001}
T.~Jacobson and D.~Mattingly (2001).
\newblock Gravity with a dynamical preferred frame.
\newblock \emph{Physical Review D} \textbf{64}, 024028.
\newblock \doi{10.1103/PhysRevD.64.024028}

\bibitem[Deser(1970)]{deser1970}
S.~Deser (1970).
\newblock Self-interaction and gauge invariance.
\newblock \emph{General Relativity and Gravitation} \textbf{1}, 9--18.
\newblock \doi{10.1007/BF00759198}

\bibitem[Jacobson(2008)]{jacobson2008}
T.~Jacobson (2008).
\newblock Einstein--\Aether\ gravity: a status report.
\newblock \emph{PoS(QG-PH)}, 020.
\newblock \doi{10.22323/1.043.0020}

\bibitem[Foster \& Jacobson(2006)]{foster2006}
B.~Z. Foster and T.~Jacobson (2006).
\newblock Post-Newtonian parameters and constraints on Einstein--\Aether\
  theory.
\newblock \emph{Physical Review D} \textbf{73}, 064015.
\newblock \doi{10.1103/PhysRevD.73.064015}
\newblock [\emph{New --- added to support explicit PPN expressions
  \eqref{eq:alpha1}--\eqref{eq:alpha2}.}]

\bibitem[Abbott et al.(2017a)]{abbott2017gw}
B.~P. Abbott et al.\ (LIGO/Virgo) (2017).
\newblock GW170817: Observation of Gravitational Waves from a Binary Neutron
  Star Inspiral.
\newblock \emph{Physical Review Letters} \textbf{119}, 161101.
\newblock \doi{10.1103/PhysRevLett.119.161101}

\bibitem[Abbott et al.(2017b)]{abbott2017mm}
B.~P. Abbott et al.\ (LIGO/Virgo) (2017).
\newblock Multi-messenger Observations of a Binary Neutron Star Merger.
\newblock \emph{Astrophysical Journal Letters} \textbf{848}, L12.
\newblock \doi{10.3847/2041-8213/aa91c9}

\end{thebibliography}



\end{document}